\documentclass[10pt,a4paper]{amsart}
\usepackage{amsmath,amssymb,amsthm,graphicx}
\usepackage{verbatim}
\newenvironment{reference}{\comment}{\endcomment}
%\newenvironment{reference}{}{}
\newenvironment{slogan}{\comment}{\endcomment}
\newenvironment{history}{\comment}{\endcomment}

% For commutative diagrams we use Xy-pic
\usepackage[all]{xy}

% We use 2cell for 2-commutative diagrams.
\xyoption{2cell}
% Theorem environments.
%
\theoremstyle{plain}
\newtheorem{theorem}[subsection]{Theorem}
\newtheorem{proposition}[subsection]{Proposition}
\newtheorem{lemma}[subsection]{Lemma}

\theoremstyle{definition}
\newtheorem{definition}[subsection]{Definition}
\newtheorem{example}[subsection]{Example}
\newtheorem{exercise}[subsection]{Exercise}
\newtheorem{situation}[subsection]{Situation}

\theoremstyle{remark}
\newtheorem{remark}[subsection]{Remark}
\newtheorem{remarks}[subsection]{Remarks}

\numberwithin{equation}{subsection}

% Macros
%
\def\lim{\mathop{\mathrm{lim}}\nolimits}
\def\colim{\mathop{\mathrm{colim}}\nolimits}
\def\Spec{\mathop{\mathrm{Spec}}}
\def\Hom{\mathop{\mathrm{Hom}}\nolimits}
\def\Ext{\mathop{\mathrm{Ext}}\nolimits}
\def\SheafHom{\mathop{\mathcal{H}\!\mathit{om}}\nolimits}
\def\SheafExt{\mathop{\mathcal{E}\!\mathit{xt}}\nolimits}
\def\Sch{\mathit{Sch}}
\def\Mor{\mathop{\mathrm{Mor}}\nolimits}
\def\Ob{\mathop{\mathrm{Ob}}\nolimits}
\def\Sh{\mathop{\mathit{Sh}}\nolimits}
\def\NL{\mathop{N\!L}\nolimits}
\def\CH{\mathop{\mathrm{CH}}\nolimits}
\def\proetale{{pro\text{-}\acute{e}tale}}
\def\etale{{\acute{e}tale}}
\def\QCoh{\mathit{QCoh}}
\def\Ker{\mathop{\mathrm{Ker}}}
\def\Im{\mathop{\mathrm{Im}}}
\def\Coker{\mathop{\mathrm{Coker}}}
\def\Coim{\mathop{\mathrm{Coim}}}

% Boxtimes
%
\DeclareMathSymbol{\boxtimes}{\mathbin}{AMSa}{"02}

%
% Macros for moduli stacks/spaces
%
\def\QCohstack{\mathcal{QC}\!\mathit{oh}}
\def\Cohstack{\mathcal{C}\!\mathit{oh}}
\def\Spacesstack{\mathcal{S}\!\mathit{paces}}
\def\Quotfunctor{\mathrm{Quot}}
\def\Hilbfunctor{\mathrm{Hilb}}
\def\Curvesstack{\mathcal{C}\!\mathit{urves}}
\def\Polarizedstack{\mathcal{P}\!\mathit{olarized}}
\def\Complexesstack{\mathcal{C}\!\mathit{omplexes}}
% \Pic is the operator that assigns to X its picard group, usage \Pic(X)
% \Picardstack_{X/B} denotes the Picard stack of X over B
% \Picardfunctor_{X/B} denotes the Picard functor of X over B
\def\Pic{\mathop{\mathrm{Pic}}\nolimits}
\def\Picardstack{\mathcal{P}\!\mathit{ic}}
\def\Picardfunctor{\mathrm{Pic}}
\def\Deformationcategory{\mathcal{D}\!\mathit{ef}}
\title{A free complex mapping to a module maps into any resolution, and any two such maps are homotopic}
\author{The Stacks Project Authors}
\date{stacks-project:algebra.tex:local-snapshot-20261012; GNU FDL 1.2 or later}
\begin{document}
\maketitle
\noindent\textit{Editorial note.} All delivered mathematical text is the author's own, copied from the pinned Stacks Project chapter file algebra.tex; no statement or proof step was rewritten, reordered or removed.\par\medskip
\noindent\textit{Editorial note.} Preamble provenance, as required for this source: the frozen snapshot stages the chapter's own preamble file beside it (preamble.tex, 217 lines, 8116 bytes, the same shared Stacks-wide preamble the published chow and etale-cohomology cards use); the first line of the pinned chapter file is \textbackslash{}input\{preamble\}, so no preamble text lives inside the chapter itself. The card takes that file's lines 8-18 (the verbatim package and the three \textbackslash{}comment-based environments reference/slogan/history at lines 9-12, section 13.2, together with the Xy-pic setup) and its lines 154 to the runtime end of the file (the native theorem declarations and mathematical macros). Not one character was changed, no package was added, no macro was defined, and the excerpt serves only as preamble, never as a span coordinate.\par\medskip
\noindent\textit{Editorial note.} The same-chapter pointers are located at their \textbackslash{}label lines inside the pinned chapter file; the cross-chapter pointers are located in the project's own tag list staged in the snapshot (sources/stacks-chow/tags.txt), at the verbatim lines carrying their canonical tags.\par\medskip
\noindent\textit{Editorial note.} The counted statement contains an Xy-pic square; the preamble spans of this card run 8-18 of the chapter's own preamble file and therefore include the Xy-pic package and option setup at lines 15-18.\par\medskip
\noindent\textit{Editorial note.} The counted spans contain no \textbackslash{}ref and no \textbackslash{}cite, so the card carries no navigation stub and no bibliography.\par\medskip
\noindent\textit{Allowed standard prerequisites (editorial).}\begin{itemize}
\item Lifting along a surjection out of a free module. Used repeatedly in the counted proof (lines 17740-17749). Produces the comparison maps degree by degree from freeness of the complex terms.
\item Exactness of a resolution for lifting into kernels. Used in both parts of the counted proof (lines 17746 and 17759). Guarantees that whatever maps to zero can be lifted one degree further.
\end{itemize}\medskip
\begin{lemma}
\label{lemma-compare-resolutions}
Let $R$ be a ring. Let $M \to N$ be a map of $R$-modules.
Let $N_\bullet \to N$ be an arbitrary resolution.
Let
$$
\ldots \to F_2 \to F_1 \to F_0 \to M
$$
be a complex of $R$-modules where each $F_i$ is a free $R$-module. Then
\begin{enumerate}
\item there exists a map of complexes $F_\bullet \to N_\bullet$ such that
$$
\xymatrix{
F_0 \ar[r] \ar[d] & M \ar[d] \\
N_0 \ar[r] & N
}
$$
is commutative, and
\item any two maps $\alpha, \beta : F_\bullet \to N_\bullet$ as in (1)
are homotopic.
\end{enumerate}
\end{lemma}

\begin{proof}
Proof of (1). Because $F_0$ is free we can find a map $F_0 \to N_0$
lifting the map $F_0 \to M \to N$. We obtain an induced
map $F_1 \to F_0 \to N_0$ which ends up in the image
of $N_1 \to N_0$. Since $F_1$ is free we may lift this
to a map $F_1 \to N_1$. This in turn induces a map
$F_2 \to F_1 \to N_1$ which maps to zero into
$N_0$. Since $N_\bullet$ is exact we see that
the image of this map is contained in the image
of $N_2 \to N_1$. Hence we may lift to get a map
$F_2 \to N_2$. Repeat.

\medskip\noindent
Proof of (2). To show that $\alpha, \beta$ are homotopic it suffices
to show the difference $\gamma = \alpha - \beta$ is homotopic
to zero. Note that the image of $\gamma_0 : F_0 \to N_0$
is contained in the image of $N_1 \to N_0$. Hence we may lift
$\gamma_0$ to a map $h_0 : F_0 \to N_1$. Consider the map
$\gamma_1' = \gamma_1 - h_0 \circ d_{F, 1}$. By our choice of $h_0$
we see that the image of $\gamma_1'$ is contained in
the kernel of $N_1 \to N_0$. Since $N_\bullet$ is exact
we may lift $\gamma_1'$ to a map $h_1 : F_1 \to N_2$.
At this point we have $\gamma_1 = h_0 \circ d_{F, 1}
+ d_{N, 2} \circ h_1$. Repeat.
\end{proof}

\end{document}
